\section{Sintaxis de coordenadas}
\label{subsec:sintaxis-de-coordenadas}

Algo que distingue a \pgfgoname\ de otros paquetes es la posibilidad de introducir coordenadas con diversas sintaxis, ofreciendo flexibilidad para representar posiciones en el tablero de manera intuitiva. 
\begin{itemize}
	\item Posibilidad de tener múltiples entradas separando con comas.
	\item Posibilidad de eliminar piedras agregando el grupo sustractivo despu\'es de un <<;>>.
\end{itemize}

Estas formas solo funcionan en los comandos plurales (|\stones|, |\marks|, |\prisoners|); |\stone|, |\mark| y |\prisoner| reciben una sola intersección. |\moves| acepta solo las que tienen un orden claro (pág.~\pageref{subsec:movimientos}). Sin embargo, por motivos de rendimiento y consistencia, se han restringido algunas libertades que en un comienzo no estaban. Por ejemplo, para referirse a una intersección, la coordenada alfabética debe escribirse en mayúsculas y sin espacios. Así, la intersección B5 debe ingresarse como |B5|, y no como |B 5|, |B  5| o |b5|, ya que los espacios entre letra y número no son válidos en este contexto. Las restricciones se explicarán en cada caso.


Cuando \pgfgoname\ tiene que leer el argumento de las coordenadas se puede encontrar con algo como esto
\begin{center}
	|\stones{|{\color{red}|<item>, <item>, |\dots;} {\color{blue} |<item>, <item>, |\dots; }{\color{red} \dots}|}|
\end{center}

Lo primero que intentar\'a hacer es dividir el argumento entre los |;| y construir dos grupos: el grupo aditivo y el grupo sustractivo. El grupo aditivo es el encargado de a\~nadir piedras mediante nuestra sintaxis. El grupo sustractivo es el encargado de eliminarlas. Estos grupos alternan automáticamente al colocar un |;|. Esto se ve con mas detalle en la subsecci\'on correspondiente. Si no se coloca un |;|, \pgfgoname\ a\~nadir\'a todas las piedras al grupo aditivo. Se pueden a\~nadir varias piedras o varios grupos separados por |,|. En conjunto, todo esto provee una interfaz para crear grupos complejos con pocos \emph{tokens};

\begin{codeexample}[width=0pt]
	\pgfgoset{profile = pgfgosmooth}
	
	\begin{goban}[17, partial = A1:Q5, label = {at = {north, west}}]
		\stones[b]{*;B,D,H,L,Q2,P,F2:4,J2rK5,N[1:2,4];B3}
	\end{goban}
\end{codeexample}


\subsection{Grupos de piedras}

\def\sssseccion#1{\vskip0.5em\noindent\textbf{#1}\par}

\sssseccion{Todas las piedras del tablero (\emph{?`Por qu\'e no?})}

A\~nade al correspondiente grupo todas las intersecciones \textbf{visibles} del tablero. Es decir, si  es un tablero parcial, solo se a\~nadir\'an las correspondientes a la regi\'on visible. Es importante aclarar esto por si se usa en combinaci\'on con |resume|. Su sintaxis es |*|.

\begin{codeexample}[]
	\pgfgoset{goban = {5, partial = B2:D4}}
	
	\begin{goban}[label = {at = {south, west}}]
		\stones[b]{*}
	\end{goban}
\end{codeexample}


\sssseccion{Intersección}

Las intersecciones se especifican indicando la letra seguida del número, sin espacios entre el componente alfabético y el numérico: |AN|

\begin{codeexample}[]
	\pgfgoset{goban = {5, partial = B2:D4}}
	
	\begin{goban}[label = {at = {south, west}}]
		\stones[b]{C2, D3}
		\stone[w]{B4}
		\stone[w]{B 3} % No se reconoce porque tiene un espacio.
	\end{goban}
\end{codeexample}

En el ejemplo, |B 3| no se reconoce debido al espacio, mientras que |B4| se dibuja correctamente como una piedra blanca.

\sssseccion{Rectas verticales u horizontales}
Las líneas son grupos de piedras verticales u horizontales delimitadas impl\'icitamente por el tama\~no del tablero: una columna o una fila entera.

\begin{itemize}
	\item Para una columna (una línea vertical), la sintaxis es |L|: |D| es la columna D.
	\item Para una fila (una línea horizontal), la sintaxis es |N|: |4| es la fila 4.
\end{itemize}

\begin{codeexample}[width=6cm]
	\pgfgoset{goban = {5, partial = B2:D4}}
	
	\begin{goban}
		\stones[b]{B2,B3,B4}
		\stones[fill = {color = red}]{D}
	\end{goban}
	
	\begin{goban}
		\stones[b]{B2,C2,D2}
		\stones[fill = {color = blue}]{4}
	\end{goban}
	
	\begin{goban}
		%Agrego las líneas C y 3. Luego elimino C3.
		\stones[fill = {color = orange}]{C, 3; C3} %OJO C,3 != C3.
	\end{goban}
\end{codeexample}

\sssseccion{Segmentos}

Los segmentos son grupos de piedras verticales u horizontales explícitamente delimitados. Como son unidireccionales (solo verticales o solo horizontales), una de sus coordenadas permanece constante y no necesita repetirse; solo se especifica el rango de la otra variable.

\begin{itemize}
	\item Para segmentos horizontales, la sintaxis es |L:LN| o |L:L N|.
	\item Para segmentos verticales, la sintaxis es |LN:N| o |L N:N|.
\end{itemize}

\begin{codeexample}[width=6cm]
	\pgfgoset{goban = {5, partial = B2:D4}}
	
	\begin{goban}
		\stones[b]{B2,B3,B4}
		\stones[w]{C2,D2}
	\end{goban}
	
	\begin{goban}
		\stones[b]{B2:4}
		\stones[w]{C:D2}
	\end{goban}
	
	\begin{goban}
		\stones[b]{B 2:4} % Los espacios funcionan entre elementos distintos.
		\stones[w]{C:D 2}
	\end{goban}
\end{codeexample}

Aquí, |B2:4| representa un segmento vertical de B2 a B4, mientras que |C:D2| indica un segmento horizontal de C2 a D2.

\begin{codeexample}[width=6cm]
	\begin{goban}[label = {at = all}]
		\stones[fill = {color = red}]{A:E5}
		\stones[fill = {color = blue}]{H9:2}
		\stones[fill = {color = orange}]{E:A2}
		\stones[fill = {color = magenta}]{A:D7, A:C9} % Pueden combinarse.
	\end{goban}
\end{codeexample}

\sssseccion{Rectángulo: una consecuencia artificial}

La sintaxis de segmentos permite crear rangos entre coordenadas para insertar piedras. Me pareció útil extenderla para incluir dos rangos y así ofrecer una primera aproximación a los rectángulos (aunque no es mi método favorito, aquí está). En este caso, la sintaxis es |L:L N:N| o |L:LN:N|. 

\begin{codeexample}[width=6cm]
	\begin{goban}[label = {at = all}]
		\stones[fill = {color = red}]{A:E 7:8} % El espacio no es obligatorio.
		\stones[fill = {color = blue}]{I:G 9:2}
		\stones[fill = {color = orange}]{A:B 2:4, F:D 1:3}
	\end{goban}
\end{codeexample}

Por ejemplo, |A:E 7:8| crea un rectángulo que abarca las columnas de A a E en las filas 7 y 8, mientras que |I:G 9:2| desciende de I9 a G2. Aunque funcional, esta sintaxis puede resultar menos intuitiva, lo que lleva a los operadores más específicos que veremos más adelante.


\sssseccion{Productos Cartesianos}
	En este contexto, el producto cartesiano es una extensi\'on m\'as poderosa de los segmentos, pero a su vez, una generalizaci\'on a todas las notaciones antes presentadas. La idea detr\'as de esto es la siguiente: Supongamos que queremos ingresar un grupo de piedras sobre la columna A: |A1:3, A5, A7:8|. Hasta el momento no hab\'ia forma mas sencilla de escribir esto. Pero todas repiten A al comienzo: ?`qu\'e tal si en lugar de escribirla tres veces la escribimos solo una? Algo como
	\begin{center}
		|A[1:3, 5, 7:8]|
	\end{center}
	En este caso, el interprete va a distribuir el A de la izquierda para cada uno de los elementos de la derecha.
	
	Los elementos de la lista de la derecha no necesariamente tienen que estar formados por n\'umeros. Si se tiene algo como esto |A:D 5:7, A1:2, A3, A:G9| tambi\'en se puede reescribir as\'i
	\begin{center}
		|A[:D 5:7, 1:2, 3, :G9]|
	\end{center}
	Claro est\'a que nadie abstrae tanto al momento de poner coordenadas, por lo que esto es solo mero capricho. La \'unica restricci\'on para el contenido entre corchetes es que no pueden anidarse corchetes. 
	
	Del mismo modo que hay una sintaxis para letras a la izquierda, hay una para n\'umeros a la derecha, y funciona igual. La expresi\'on |A12, C12, E:G12| puede escribirse como |[A,C,E:G]12|.
	
	Existe un tercer operador, que s\'i, es m\'as un capricho que otra cosa y es el producto cartesiano sin restricciones. No me detendr\'e mucho en \'el, porque este requiere mayor grado de abstracci\'on y realmente no tiene utilidad humana. Este operador es una conjunci\'on de los dos anteriores: La expresi\'on |[A,D:F][2,5,7:9]| se interpreta como
	\begin{center}
		|A2, A5, A7:9, D:F2, D:F5, D:F 7:9|.
	\end{center}
 
	\begin{codeexample}[width=4cm]
		\begin{goban}
			\stones[b]{[A, D:F][2, 5, 7:9]}
			\marks[b]{A2, A5, A7:9, D:F2, D:F5, D:F 7:9} %<- lo que lee el interprete.
		\end{goban}
	\end{codeexample}

	La sintaxis de producto cartesiano entonces se puede invocar de tres formas distintas:
\begin{itemize}
	\item |L[itm1, itm2,...]|.
	\item |[itm1, itm2,...]N|.
	\item |[itm1,itm2,...][itm3,itm4,...]|.
\end{itemize}
donde los corchetes no pueden anidarse. 

Esto puede combinarse tambi\'en con algunos operadores que veremos a continuaci\'on para crear nodos en los caminos.

Hasta aquí hemos visto comandos para intersecciones que no pueden combinarse con otros en una misma instrucción y que, por lo tanto, \textbf{deben} separarse por comas. Ahora exploraremos comandos que sí permiten encadenarse entre sí, ofreciendo mayor flexibilidad para patrones complejos.

\sssseccion{Horizontal--Vertical y Vertical--Horizontal}

Dadas dos intersecciones no alineadas vertical ni horizontalmente, estos comandos generan las intersecciones que conectan la primera con la segunda mediante un recorrido «horizontal--vertical» o «vertical--horizontal». Esto puede recordar al operador \verb+-|+ de TikZ, de donde tomé prestada la idea. La sintaxis es \verb+LN -| LN+ (para el primer caso) o \verb+LN |- LN+ (para el segundo), con o sin espacios. Estos comandos, junto con el siguiente, pueden encadenarse para formar escaleras o recorridos propios de la m\'etrica de Manhattan.

\begin{codeexample}[width=6cm]
	\begin{goban}[label = {at = all}]
		\stones[fill = {color = red}]{A1 -| D3}
		\stones[fill = {color = blue}]{D4 -| A2}
		\stones[fill = {color = orange}] {A9 -| C7 -| G5} % Encadenados.
		\stones[fill = {color = magenta}] {F1 |- I3 -| E4} % Intente adivinar.
	\end{goban}
\end{codeexample}

En este ejemplo, \verb+A1 -| D3+ va horizontalmente de A1 a D1 y luego verticalmente a D3.

\sssseccion{Línea horizontal o vertical}

Dadas dos intersecciones alineadas vertical u horizontalmente, este operador las conecta con una línea recta de piedras. Su sintaxis es |LN -- LN|. Puede encadenarse con los dos comandos anteriores para formar patrones más elaborados. Si las intersecciones no están alineadas, el operador no dibuja nada: para eso están las líneas en cualquier dirección y las escaleras, más abajo.

\begin{codeexample}[width=6cm]
	\begin{goban}[label = {at = all}]
		\stones[fill = {color = red}]{A1 -- A5}
		\stones[fill = {color = blue}] {A1 -- D1}
		\stones[fill = {color = magenta}] {A6 -- B9} % No dibuja porque no están alineados.
		\stones[fill = {color = orange}] {I9 -| F6 -- I6} % Encadenado.
	\end{goban}
\end{codeexample}

Aquí, |A1 -- A5| dibuja una línea vertical, mientras que \verb+I9 -| F6 -- I6+ combina un recorrido «horizontal--vertical» con una línea horizontal adicional.

¡Ojo! \verb+A1 -- D1 -| F3+ es un encadenamiento permitido, mientras que \verb+A:D1 -| F3+ no lo es, porque el |:| de los rangos no es un operador de cadena.

\sssseccion{Líneas en cualquier dirección y escaleras}

Para dos intersecciones cualesquiera, |LN :--: LN| dibuja la línea de piedras más parecida a la recta que las une: una piedra por fila (o por columna, si la línea avanza más de lo que sube o baja), la más cercana a la recta. Esas piedras se tocan solo por las esquinas en cada escalón, como en una diagonal, así que en go no forman un grupo. Para que lo formen, \verb+LN :-|: LN+ y \verb+LN :|-: LN+ rellenan cada escalón con la esquina de \verb+-|+ (primero horizontal) o de \verb+|-+ (primero vertical): el resultado es una escalera conectada. Si las dos intersecciones están en una misma fila o columna, los tres operadores dan la misma línea recta.

\begin{codeexample}[width=6cm]
	\begin{goban}[label = {at = all}]
		\stones[fill = {color = red}]{A1 :--: C7}
		\stones[fill = {color = blue}]{D1 :-|: F7}
		\stones[fill = {color = orange}]{G1 :|-: I7}
	\end{goban}
\end{codeexample}

Cuando la recta pasa justo entre dos intersecciones, la piedra va del lado del extremo más cercano, y en el medio exacto, a la columna (o la fila) menor. Así, |A1 :--: C5| y |C5 :--: A1| dan las mismas piedras, y la línea queda simétrica, como la recta. En las escaleras el orden sí importa, como en \verb+-|+ y \verb+|-+.

Las intersecciones van en el orden del camino, de la primera a la segunda, con la esquina de cada escalón antes del escalón: |\moves| y |sequence| las numeran a lo largo de la línea. Los tres operadores se encadenan con los demás, como en \verb+A1 :--: E3 -| G6+; los dos puntos son parte del operador, así que no se confunden con un rango.

\begin{codeexample}[width=6cm]
	\begin{goban}[7, label = {at = {south, west}}]
		\stones[w]{A1 :-|: F4}
		\marks[sequence]{A1 :-|: F4}
	\end{goban}
\end{codeexample}

\sssseccion{Rectángulos rellenos y vacíos}

La sintaxis previa para rectángulos |L:L N:N| es poco intuitiva y más bien una herencia de los segmentos acotados. Por ello, existen dos operadores más directos para crear rectángulos: |R| y |r|. La diferencia radica en que |R| genera un rectángulo relleno, mientras que |r| solo dibuja su contorno. Dados dos vértices diametralmente opuestos del rectángulo, la sintaxis es |LN R LN| o |LN r LN|. Aunque los espacios no son obligatorios por ahora, recomiendo usarlos para mayor legibilidad.

\begin{codeexample}[width=6cm]
	\begin{goban}[label = {at = all}]
		\stones[fill = {color = red}]{A1 R D3}
		\stones[fill = {color = blue}] {A5 r D8}
		\stones[fill = {color = orange}] {I8 r G4, H6} % ¡Está vivo!
	\end{goban}
\end{codeexample}

Por ejemplo, |A1 R D3| rellena un rectángulo de A1 a D3, mientras que |A5 r D8| dibuja solo el borde de un rectángulo mayor.

Los rectángulos también son operadores de cadena y pueden usarse como tales para crear patrones más complejos.

\begin{codeexample}[width=6cm]
	\begin{goban}[label = {at = all}]
		\stones[fill = {color = orange}] {I8 r G6 r I4} % Ud. adivine.
		\marks[sequence]{I8,G6,I4}
		\stones[fill = {color = blue!20}] {A1 R E3 r A6 |- E9 -- E8} % ¡Wow!
		\marks[sequence]{A1,E3,A6,E9,E8}
	\end{goban}
\end{codeexample}

En este caso, |I8 r G6 r I4| encadena dos rectángulos de contorno, mientras que \verb+A1 R E3 r A6 |- E9 -- E8+ combina un rectángulo relleno, un contorno, un recorrido «vertical--horizontal» y una línea, mostrando la potencia del encadenamiento.

\subsection{Grupos aditivos y sustractivos}

Supongamos que queremos dibujar el siguiente diagrama:

\begin{figure}[H]
	\centering
	\begin{goban}
		\stones[b]{A1 R E6; C6 |- E3, A1:2;}
	\end{goban}
\end{figure}

Esto puede lograrse con los operadores ya vistos, dividiendo el patrón en varias partes:

\begin{codeexample}[width=6cm]
	\begin{goban}[label = {at = all}]
		\stones[fill = {color = orange}] {A6 R B3, B2 R E1, D4 R E6}
	\end{goban}
\end{codeexample}

Pero también podemos imaginar que partimos de un rectángulo completo al que le hemos quitado las piedras marcadas, ofreciendo una perspectiva más conceptual:

\begin{figure}[H]
	\centering
	\begin{goban}
		\stones[b]{A1 R E6}
		\marks[cross, line = {color = white}]{C6 |- E3, A1:2}
	\end{goban}
\end{figure}

De esto tratan los grupos aditivos y sustractivos. Un grupo aditivo añade intersecciones a la lista, mientras que uno sustractivo las elimina. Por defecto, los comandos siempre comienzan procesando grupos aditivos, lo que explica que dibujen piedras directamente. En \pgfgoname, podemos alternar entre ambos tipos usando |;|.

\begin{codeexample}[width=6cm]
	\begin{goban}[label = {at = all}, partial = A1:E6]
		\stones[b]{A1 R E6} % Grupo aditivo.
	\end{goban}
	
	\begin{goban}[label = {at = all}, partial = A1:E6]
		\stones[b]{A1 R E6; A 1:2} % Grupo aditivo; sustractivo.
	\end{goban}
	
	\begin{goban}[label = {at = all}, partial = A1:E6]
		\stones[b]{A1 R E6; A 1:2, C6 |- E3} % Grupo aditivo; sustractivo.
	\end{goban}
\end{codeexample}

En esta secuencia, primero se dibuja un rectángulo completo (|A1 R E6|), luego se elimina el segmento horizontal |A1:2|, y finalmente se quita el recorrido \verb+C6 |- E3+. Esto resulta en el mismo diagrama que el ejemplo inicial, pero con un enfoque más compacto.

La ventaja de este sistema es que permite realizar ajustes cuando sea necesario. A veces, un patrón aditivo es ideal salvo por una o dos piedras; en esos casos, podemos eliminarlas con un grupo sustractivo sin rehacer todo el diseño.

\begin{codeexample}[width=6cm]
	\begin{goban}[label = {at = all}]
		\stones[b]{A1 R E6; B:D 5, C2}
	\end{goban}
\end{codeexample}

Aquí, |B:D 5| elimina un segmento horizontal en la fila 5, y |C2| quita una piedra específica, ajustando el rectángulo inicial.

Recordemos que los grupos aditivos y sustractivos pueden encadenarse, separándolos con |;|, lo que permite alternar entre añadir y quitar elementos en una sola línea tantas veces como se desee.

\begin{codeexample}[width=6cm]
	\begin{goban}[label = {at = all}, partial = A1:E6]
		\stones[b]{A1 R E6} % AD
	\end{goban}
	
	\begin{goban}[label = {at = all}, partial = A1:E6]
		\stones[b]{A1 R E6; A1 R C4, E6} % AD, SU
	\end{goban}
	
	\begin{goban}[label = {at = all}, partial = A1:E6]
		\stones[b]{A1 R E6; A1 R C4, E6; A 1:2} % AD, SU, AD
	\end{goban}
\end{codeexample}

En este último ejemplo, se crea un rectángulo, se eliminan dos regiones (|A1 R C4| y |E6|), y luego se añade un segmento (|A 1:2|).
